% Part: second-order-logic
% Chapter: sol-and-set-theory

\documentclass[../../../include/open-logic-chapter]{subfiles}

\begin{document}

\olchapter{sol}{set}{દ્વિતીય-ક્રમ તર્કશાસ્ત્ર અને ગણસિદ્ધાંત}

\begin{editorial}
આ પ્રકરણ દ્વિતીય-ક્રમ તર્કશાસ્ત્રમાં ઘાતગણો અને સાતત્યકનું
સંકેતીકરણ કરે છે. સ્થિર અંગ્રેજી મૂળમાં અનેક પરિણામોની સાબિતીઓ
ઉમેરવાની બાકી છે, અભ્યાસપ્રશ્નો નથી અને વ્યાખ્યાઓ તથા પરિણામો
ભૂલવાળાં હોઈ શકે એવી સંપાદકીય નોંધ છે. આ ગુજરાતી આવૃત્તિમાં
નિર્દિષ્ટ સ્રોત-સુધારાઓ અને કેટલીક પૂરક દલીલો આપી છે; મૂળના
અધૂરા પ્રકરણ અંગેની નોંધ જાળવી છે. સાવચેતીથી વાંચો અને કંઈ
ખોટું કે અસ્પષ્ટ લાગે તો તેની જાણ કરો.
\end{editorial}

\olimport{introduction}

\olimport{comparing-sets}

\olimport{cardinalities}

\olimport{power-of-continuum}

\OLEndChapterHook

\end{document}
